% Proof-ready supplement for Part B. The parent may renumber B.14a.
\begin{statement}{B.14a}{Proved}{An integral group lift of the exceptional-prime witness through class thirteen}
Let $F=\langle x,y\rangle$ be the discrete free group,
$N_{13}=F/\gamma_{14}F$, and $z=(xy)^{-1}$.
There is an element $f\in\gamma_7N_{13}$ whose image in
$\gamma_7F/\gamma_8F$ is the specific integral vector
$q_7$ of B.12 and which satisfies the genuine group equations
\[
 f(y,x)f(x,y)=1,\qquad
 f(z,x)f(y,z)f(x,y)=1
 \quad\hbox{in }N_{13}.
\]
The same assertion holds with $q_7$ replaced by any integral
multiple. Consequently these finite nilpotent jets exist over
$\mathbb Z_2$ and $\mathbb Z_3$ without scaling $q_7$.
This is a finite-stage assertion; it does not assert an infinite
pro-$2$, pro-$3$, or full profinite lift.
\end{statement}
\begin{proof}
We use actual group commutators $[g,h]=ghg^{-1}h^{-1}$.
For a Lyndon word $w$, let $C_x=x$, $C_y=y$, and
$C_w=[C_u,C_v]$, where $w=uv$ is the standard Lyndon
factorization with $v$ the longest proper Lyndon suffix.
The leading Magnus term of $C_w-1$ is the standard integral
Lie bracket $b_w$. These leading terms form a basis of each
$\gamma_nF/\gamma_{n+1}F$. Induction over the lower central
series therefore gives unique integral coordinates for
$\gamma_7F/\gamma_{14}F$ in the ordered products of
$C_w$, $7\leq |w|\leq13$.

Set $x=1+X$, $y=1+Y$ in the truncated noncommutative
power-series ring over $\mathbb Z$. In particular
\[
 z=(1+Y)^{-1}(1+X)^{-1},
\]
and this full group series is used in every substitution.
If $A=g-1$, $B=h-1$, the exact identity
\[
 [g,h]-1=(AB-BA)(1+A)^{-1}(1+B)^{-1}
\]
constructs every needed Hall series by finite integer arithmetic.
The product of two terms of degree at least seven has degree at
least fourteen. Hence through degree thirteen the Magnus expansion
of $\prod C_w^{a_w}$ is exactly
$1+\sum a_w(C_w-1)$, and both displayed group equations
are linear equations in all the Hall exponents.

We fix the seven weight-seven exponents to the coefficients of
$q_7$ and include \emph{every} exponent in weights eight through
thirteen as an unknown. The simultaneous systems, retaining all
earlier free parameters, have the following exact ranks:
\[
\begin{array}{c|rrrrrr}
\text{class}&8&9&10&11&12&13\\\hline
\text{unknown exponents}&30&86&185&371&706&1336\\
\text{rank}&20&66&147&302&578&1103\\
\text{free parameters}&10&20&38&69&128&233.
\end{array}
\]
The computations through class twelve are recorded in
\texttt{code/e\_exceptional\_jets.py} and its class-twelve JSON
output. The final class-thirteen calculation uses
\path{code/e_exceptional_lattice.py} and the exact certificate
\texttt{e\_exceptional\_lattice\_13.json}.
The complete sparse matrix, its independent row indices, and the
right-hand side are also retained in the compressed JSON file
\texttt{e\_exceptional\_matrix\_13.json.gz}.

Here is the integral certificate extraction, so that existence
is not inferred from a rational solution. Write the final system
as $Aa=b$, with $2708$ rows and $1336$ unknowns. Select $1103$
independent rows, producing $Ma=y$. The rank over $\mathbb Q$
is checked exactly and agrees with the selected rank. Let $H_0$
be the nonsingular block of the row Hermite normal form of
$M^{\mathsf t}$. Its rows are a basis of the integer row lattice,
so the columns of $H_0^{\mathsf t}$ are a basis of
$M\mathbb Z^{1336}$. The computation verifies
\[
 z=(H_0^{\mathsf t})^{-1}y\in\mathbb Z^{1103},
 \qquad
 R=(H_0^{\mathsf t})^{-1}M
       \in M_{1103,1336}(\mathbb Z).
\]
In particular the minimum integral multiplier of the specified
leading vector is one. The columns of $R$ generate the whole
lattice $\mathbb Z^{1103}$.

For explicit coordinates, choose full-rank column minors of $R$
modulo $2$ and modulo $3$. Inverting these minors over $\mathbb Q$
gives solutions $a_p=U_p/d_p$, $p=2,3$, supported on the chosen
columns, where $U_p$ is integral and $d_p$ is prime to $p$.
Every equation of the original system $Aa_p=b$ is checked.
The two computed denominators are coprime. The JSON certificate
records the denominators, selected columns and exact B\'ezout
coefficients $s_2,s_3$ satisfying
\[
 s_2d_2+s_3d_3=1.
\]
Consequently
\[
 a=s_2U_2+s_3U_3\in\mathbb Z^{1336},\qquad Aa=b.
\]
The certificate's \texttt{hall\_coordinates} field records this
integral vector together with the seven fixed coefficients of
$q_7$, and therefore specifies the group word completely.
The large integer coefficients are kept in that machine-readable
certificate.

The final verification reconstructs the Hall commutators and
checks \emph{every associative-word coefficient} of both group
defects through degree thirteen, with exact integer arithmetic.
Both reconstructed defects are zero. Vanishing in the truncated
Magnus ring implies the asserted identities in $N_{13}$: if a
group word had a first nonzero lower-central component of degree
at most thirteen, its nonzero Lyndon Lie polynomial would be its
first nonzero Magnus term. Conversely, a word in $\gamma_{14}F$
has no such term.

For completeness, selecting Lyndon-word rows when solving does
not discard an equation. Each cycle defect is primitive, up to
the given truncation, for the Magnus Hopf structure
$\Delta X=X\otimes1+1\otimes X+X\otimes X$.
Its first nonzero homogeneous part would be an ordinary free
Lie polynomial. The triangular Lyndon basis detects that part
on Lyndon words. Induction on degree proves that vanishing of
all homogeneous Lyndon coefficients is equivalent to vanishing
of the entire defect. The independent reconstruction of the full
associative expansions provides a further direct verification
of the recorded solution.

Finally, in this degree range replacing every Hall exponent by
its multiple $m$ multiplies both linear defects by $m$ and
replaces the leading term by $m q_7$. This proves the assertion
for all integral multiples, including the two multiples singled
out in the prompt. Passing to the pro-$p$ completion of $N_{13}$
or any finite $p$-group quotient preserves the verified equations.
\end{proof}

\paragraph{What remains of the exceptional-prime problem.}
The calculation evaluates the actual group-equation lifting
problem through the stated nilpotency class. It does not prove
that one of these jets extends through every subsequent class.
Existence at all classes, with integral $p$-adic coordinates,
would suffice: compactness of the product of the Hall-coordinate
copies of $\mathbb Z_p$ would give a compatible inverse-limit
solution. The required nonempty sets have been checked only to
the finite bound above. In particular the factor $6$ in an
averaging argument is not a unit for the two primes in question;
taking norms may multiply the first term at each new stage and
therefore does not establish an infinite lift of $q_7$, $2q_7$,
or $3q_7$.
